\section{Algorithm}
\label{sec:algorithm}

\subsection{Ionizing photon rate: blackbody spectrum fit}
\label{sec:blackbody}

Every active star's ionizing photon rate $\dot N_{\rm ion}$ (the budget
spent in Section~\ref{sec:photon-budget} below) comes from
\code{radiation\_get\_individual\_star\_ionizing\_photon\_emission\_
rate\_fit} (\code{src/feedback/GEAR/radiation.c}), used both directly
for individual (non-SSP) star particles and to build the IMF-integrated
lookup table population stars read from. Unlike a simple empirical
$\dot N_{\rm ion}(M)$ fit, this derives $\dot N_{\rm ion}$ from the
star's own blackbody spectrum, via a mass $\to$ radius $\to$
luminosity $\to$ effective temperature chain:
\begin{align}
  R(M) &= R_\odot \times
  \begin{cases}
    (M/M_\odot)^{0.8} & M < 1\,M_\odot \\
    (M/M_\odot)^{0.57} & 1 \le M/M_\odot < 8 \\
    (M/M_\odot)^{0.5} & M/M_\odot \ge 8
  \end{cases}
  \label{eq:mass-radius} \\
  L(M) &= L_\odot \times
  \begin{cases}
    0.185\,(M/M_\odot)^{2} & M/M_\odot < 0.43 \\
    (M/M_\odot)^{4} & 0.43 \le M/M_\odot < 2 \\
    1.5\,(M/M_\odot)^{3.5} & 2 \le M/M_\odot < 54 \\
    32000\,(M/M_\odot) & M/M_\odot \ge 54
  \end{cases}
  \label{eq:mass-luminosity} \\
  T_{\rm eff} &= 5780\,{\rm K} \times
  \left(\frac{L/L_\odot}{(R/R_\odot)^2}\right)^{1/4},
  \label{eq:stefan-boltzmann}
\end{align}
i.e.\ a piecewise power-law mass--radius and mass--luminosity relation
(\code{radiation\_get\_individual\_star\_radius},
\code{radiation\_get\_individual\_star\_luminosity}), then $T_{\rm eff}$
from the Stefan--Boltzmann law applied to $L$ and $R$ directly (not a
separate mass--temperature fit -- see the note at the end of this
subsection). $\dot N_{\rm ion}$ is then the photon-number flux of a
Planck spectrum of temperature $T_{\rm eff}$ above the H-ionizing edge
$h\nu_0=13.6\,$eV:
\begin{equation}
  \dot N_{\rm ion} = \frac{L}{k_B T_{\rm eff}} \times
  \frac{15}{\pi^4} \int_{x_0}^{\infty} \frac{x^2}{e^x-1}\,dx, \qquad
  x_0 = \frac{h\nu_0}{k_B T_{\rm eff}},
  \label{eq:planck-photon-rate}
\end{equation}
the standard Planck photon-number integral normalized so the
$15/\pi^4$ prefactor makes the \emph{bolometric} photon rate exact in
the $x_0\to0$ limit. The incomplete integral in
Eq.~\ref{eq:planck-photon-rate} has no closed form, so the code
evaluates it via the standard series expansion
\begin{equation}
  \int_{x_0}^\infty \frac{x^2}{e^x-1}\,dx =
  \sum_{n=1}^{\infty} \frac{e^{-nx_0}}{n}
  \left(x_0^2 + \frac{2x_0}{n} + \frac{2}{n^2}\right),
  \label{eq:planck-series}
\end{equation}
truncated at $n=5$ terms (or earlier once $e^{-nx_0}<10^{-10}$) --
sufficient given $x_0 \gtrsim 3$ for any star hot enough to ionize
hydrogen at all, so the series converges geometrically fast. If
$x_0 > 45$ (i.e.\ $e^{-x_0}<10^{-20}$, negligible ionizing output --
happens for any star with $T_{\rm eff}\lesssim3500\,$K), $\dot N_{\rm
ion}=0$ is returned directly rather than evaluating the series.

\textbf{Note: two independent mass--temperature relations coexist,
confirmed to disagree (investigated 2026-07-20).}
\code{radiation\_get\_individual\_star\_temperature} is a separate,
piecewise mass--temperature fit ($3500\,{\rm K}\,(M/M_\odot)^{0.5}$
below $1\,M_\odot$, $5800\,{\rm K}\,(M/M_\odot)^{0.5}$ up to
$8\,M_\odot$, $25000\,{\rm K}\,(M/M_\odot/20)^{0.1}$ above), \emph{not}
used by Eq.~\ref{eq:stefan-boltzmann} above, which derives $T_{\rm
eff}$ from $L$ and $R$ via Stefan--Boltzmann instead. A grep confirms
\code{radiation\_get\_individual\_star\_temperature} has zero call
sites anywhere in \code{src/} -- genuinely dead code, not merely
redundant. Evaluated both relations (plus the $R(M)$/$L(M)$ fits the
Stefan--Boltzmann path depends on) across $0.1$--$150\,M_\odot$: they
agree only near the two piecewise-boundary anchors ($\approx1\,M_\odot$,
$\approx8\,M_\odot$, almost certainly tuned to known reference points
like the Sun's $T_{\rm eff}$) and diverge by a factor
$\sim1.5$--$2.7\times$ everywhere else, worst in the OB-star range that
dominates ionizing photon production. The Stefan--Boltzmann path (the
one actually used) also saturates at a fixed $\approx77{,}300\,$K for
every mass above $\approx54\,M_\odot$ -- a structural artifact of the
mass--luminosity fit turning linear ($L\propto M$) at exactly the mass
the mass--radius fit is already $R\propto M^{0.5}$, making $L/R^2$ (and
hence $T_{\rm eff}$) mass-independent there by construction, not
physical intent. \textbf{Not worth reconciling in place}: both fits are
single-metallicity piecewise power laws with no $(M,Z)$ dependence at
all, which is the actual limitation regardless of which temperature
relation is used. This entire empirical-fit layer (radius, luminosity,
temperature) is intended to be replaced by extending \texttt{pychem}
(\texttt{\textasciitilde/programs/pychem/}, already the source of this
project's GEAR chemical-yield tables) to emit proper
$(M,Z)$-dependent stellar tables -- see the open items
(Section~\ref{sec:open}) for that standalone follow-up project. When it
lands, the unused \code{radiation\_get\_individual\_star\_temperature}
should simply be deleted.

\subsection{Photon budget and recombination cost}
\label{sec:photon-budget}

The whole algorithm -- one-shot, distance-sorted photon-budget spending
on the nearest candidate gas particles, closest-first, until exhausted
-- is the FIRE-2 subgrid photoionization scheme
\citep[Appendix~E]{Hopkins2018}; Hu et al.\ and Smith et al.\ (and this
branch) build directly on it. Each active star with ionizing photon
rate $\dot N_{\rm ion}$ (Eq.~\ref{eq:planck-photon-rate},
Section~\ref{sec:blackbody}) gathers candidate gas particles within a
search radius, sorted by distance, and spends its budget on them
closest-first. For a candidate particle $b$ of mass $m_b$, hydrogen
mass fraction $X_H$, and physical density $\rho_b$, the number of
hydrogen atoms and the recombination rate needed to keep it ionized
are
\begin{align}
  N_{H,b} &= \frac{X_H\, m_b}{m_p}, \\
  n_{e,b} &= \frac{X_H\, \rho_b}{m_p}, \\
  \Delta \dot N_b &= N_{H,b}\, \beta\, n_{e,b},
  \label{eq:recomb-cost}
\end{align}
where $m_p$ is the proton mass and $\beta$ is the case-B recombination
coefficient, a \emph{fixed} constant evaluated at the conventional
$T=10^4\,$K (\code{phys\_const->const\_caseb\_recomb}), matching the
convention both reference papers use. Equation~\ref{eq:recomb-cost} is
identical in form to Hu et al.'s $N_H\beta n_H$ (their $N_H$ is called
``the number of hydrogen nuclei'', an unambiguous atom count) and to
Smith et al.'s $R_{\rm rec}=\beta m \rho (X_H/m_p)^2$. Neither $N_H$
nor $n_e$ has a mean-molecular-weight ($\mu$) or metallicity factor;
see Section~\ref{sec:mu-bug} for why our first attempt at this
included one.

\textbf{Known approximation: $n_e \approx n_H$.} Equation~\ref{eq:recomb-cost}
(and every other electron-density calculation in this model, including
the adaptive rebuild cadence's recombination timescale,
Section~\ref{sec:adaptive-cadence}) assumes one free electron
per hydrogen nucleus -- it does not track helium's ionization state at all.
For a standard cosmic He mass fraction ($\sim$24--28\%, $\sim$8--10\%
by number relative to H), singly-ionized He alone (needs $>24.6\,$eV,
present in any O-star spectrum energetic enough to ionize H at all)
underestimates $n_e$ by $\sim$8--10\% if ignored; doubly-ionized He
($>54.4\,$eV, only the hottest O stars) up to $\sim$16--20\%. This is a
structural simplification of the whole model (this branch does not add
or remove it), not something any single feature here fixes -- any
future recombination-timescale or electron-density calculation added
to this model inherits the same approximation.

\subsection{Temperature dependence of $\beta$}
\label{sec:beta-temperature}

\citet{HuiGnedin1997}, Appendix A, give a fit to the case-B H\,{\sc
ii} recombination coefficient (from \citealt{Ferland1992}'s data,
accurate to 0.7\% from 1\,K to $10^9\,$K -- effectively exact for our
purposes):
\begin{equation}
  \alpha_B(T) = 2.753\times10^{-14}\,
  \frac{\lambda^{1.500}}{\left[1+(\lambda/2.740)^{0.407}\right]^{2.242}}
  \ {\rm cm^3\,s^{-1}}, \qquad
  \lambda = \frac{315614\,{\rm K}}{T}.
  \label{eq:alphaB-huignedin}
\end{equation}
This was already implemented correctly elsewhere in \swift --
identically in \code{src/rt/SPHM1RT/rt\_cooling\_rates.h}, and inside
grackle's own \code{k2\_rate()} (which GEAR-RT delegates to for its own
case-B recombination, via the \code{CaseBRecombination} flag -- GEAR-RT
has no case-B formula of its own) -- and \textbf{is now also
implemented here (2026-07-16)}: $\beta$ in Eq.~\ref{eq:recomb-cost} is
Eq.~\ref{eq:alphaB-huignedin} evaluated at $T$, not a fixed constant.
$T$ is the particle's own applied ionized temperature -- the same
$u_{\rm new}$ from Eq.~\ref{eq:floor}, converted to a temperature via
its own $\mu$ -- so the recombination cost charged against a star's
budget is always consistent with the temperature that particle is
actually being held at
(\code{radiation\_get\_part\_ionized\_internal\_energy}, shared between
\code{radiation\_get\_part\_rate\_to\_fully\_ionize} and
\code{cooling\_ionize\_part\_subgrid} to guarantee the two never
disagree; \code{radiation\_get\_case\_b\_recombination\_\allowbreak
coefficient\_cgs} implements Eq.~\ref{eq:alphaB-huignedin} itself).

For reference, evaluating Eq.~\ref{eq:alphaB-huignedin} across the
range this model's own temperature floor
(Eq.~\ref{eq:floor}--\ref{eq:tcollisional}) applies --
$\approx$8,600\,K at $Z=Z_\odot$ up to $\approx$47,500\,K at $Z=0$ --
against the value $\beta$ used to take everywhere before this fix
($2.6\times10^{-13}\,{\rm cm^3\,s^{-1}}$, $\alpha_B(10^4\,{\rm K})$ to
+0.3\%):
\begin{center}
\begin{tabular}{lrr}
  \hline
  $T$ [K] & $\alpha_B(T)$ [cm$^3$ s$^{-1}$] & vs.\ the old fixed value \\
  \hline
  8,600 ($Z=Z_\odot$)  & $2.94\times10^{-13}$ & $+13\%$ \\
  10,000 (old fixed value's own $T$) & $2.59\times10^{-13}$ & $-0.3\%$ \\
  20,000               & $1.43\times10^{-13}$ & $-45\%$ \\
  47,500 ($Z=0$)        & $6.42\times10^{-14}$ & $-75\%$ \\
  \hline
\end{tabular}
\end{center}
Since $\alpha_B(T)$ falls with rising $T$, the old fixed value used to
\emph{overestimate} the recombination cost at low metallicity by a
factor of $\sim$4 at this example's $Z=0$ default -- the star's budget
was spent roughly 4$\times$ faster than physically necessary, which
artificially shrank the achieved $r_{\rm HII}$ in exactly the direction
that could explain part of the residual gap to
\citet{Hu2017}/\citet{Smith2021}'s reference numbers
(Section~\ref{sec:metallicity}). At $Z=Z_\odot$ the effect was smaller
and opposite in sign ($+13\%$, i.e.\ a mild \emph{under}estimate). See
Section~\ref{sec:validation} for the measured effect of this fix on the
validation matrix.

A particle is a candidate for ionization only if it is not already
tagged, is denser than \code{HII\_region\_min\_density\_Hpcm3}, and
is not already above $\approx 10^4\,$K
(\code{feedback\_part\_can\_be\_ionized}). The star consumes its
budget particle by particle until it can no longer afford the next
one. The last, unaffordable particle is the ``boundary particle'':
\begin{itemize}
  \item \code{HII\_deterministic\_boundary\_ionization=1}: always
    ionize it anyway (over-spending the budget for that cycle);
  \item \code{HII\_deterministic\_boundary\_ionization=0}: ionize it
    with probability equal to the affordable fraction, unbiased in
    expectation.
\end{itemize}
See Section~\ref{sec:validation} for the measured effect of this
choice -- it turns out to matter a lot more than the name ``one
boundary particle'' suggests. \textbf{Decided (2026-07-16):
probabilistic is the default}, not deterministic -- the validation
matrix showed deterministic mode over-ionizes (mildly for the uniform
scheme, severely for the angular scheme, driving it straight into the
search-radius cap) with no compensating benefit, so this is simply the
better default rather than a papers-vs-production tradeoff. The code's
own default (\code{default\_HII\_deterministic\_boundary\_ionization},
\code{feedback\_properties.h}) was already 0; only example configs that
had explicitly overridden it needed fixing.

\subsection{Angular (HealPix) budget splitting}
\label{sec:angular-splitting}

To avoid a dense clump near a source absorbing a disproportionate
share of the photon budget regardless of the small solid angle it
subtends (``mass-biasing''), \code{HII\_angular\_nside} splits the
budget into $N_{\rm pix} = 12\,{\rm nside}^2$ equal shares
($\dot N_{\rm ion}/N_{\rm pix}$ each) and assigns each candidate
particle to exactly one pixel by its direction from the star, hard
assignment (no fractional overlap, matching both papers' own choice).
\code{nside=0} degenerates to $N_{\rm pix}=1$, i.e.\ today's plain
spherical behaviour. Only \code{nside=0} and \code{nside=1}
($N_{\rm pix}=12$) are currently supported.

\subsection{Temperature floor}

Once a particle is tagged, its internal energy is set (and reset every
step while tagged, in the cooling task, via
\code{cooling\_ionize\_part\_subgrid}) to
\begin{equation}
  u_{\rm new} = \min\left(\Delta u_{\rm ionized},\ u_{\rm collisional}\right),
  \label{eq:floor}
\end{equation}
where $\Delta u_{\rm ionized} = N_H E_{\rm ion}/m_b$ is the specific
energy needed to fully ionize the particle ($E_{\rm ion}=13.6\,$eV per
H atom, metallicity-independent), and $u_{\rm collisional}$ comes from
a metallicity-dependent equilibrium temperature. Unlike the rest of
Section~\ref{sec:algorithm}, Eq.~\ref{eq:tcollisional} is not from
FIRE-2 -- FIRE-2 clamps ionized gas to a flat $10^4\,$K -- but from
its successor, \citet[][the metallicity-dependent photoionization
temperature fit]{Hopkins2023}:
\begin{equation}
  T_{\rm collisional} = 10^4\,{\rm K} \times
  \begin{cases}
    \min\left(6.62,\ \dfrac{0.86}{1+0.22\log_{10}(Z/Z_\odot)}\right) & Z \ge 10^{-3}Z_\odot \\[2mm]
    6.62 & Z < 10^{-3}Z_\odot.
  \end{cases}
  \label{eq:tcollisional}
\end{equation}
Neither Hu et al.\ nor Smith et al.\ have a metallicity-dependent floor
at all -- both simply clamp ionized gas to a flat $10^4\,$K. See
Section~\ref{sec:metallicity} for the numeric consequence of
Eq.~\ref{eq:floor}--\ref{eq:tcollisional} at $Z=0$.

\textbf{Bug found and fixed (2026-07-18):} Eq.~\ref{eq:tcollisional} is
$\log_{10}$, but \code{radiation.c} computed it with C's \code{log()}
(natural log) until this fix. This silently applied the wrong base
everywhere $Z \ge 10^{-3}Z_\odot$, including at this project's own
$Z/Z_\odot\approx0.231$ reference point (Section~\ref{sec:metallicity}),
chosen specifically because $\log_{10}$ makes
$T_{\rm collisional}=10^4\,$K there. With the natural-log bug, the code
actually applied $T_{\rm collisional}\approx12{,}700\,$K at that same
point -- not the intended, papers-matching $10^4\,$K -- which also
shifted the case-B $\beta$ (Section~\ref{sec:beta-temperature}, evaluated
at this same temperature whenever $u_{\rm collisional} < \Delta
u_{\rm ionized}$). Every deprecated result in this logbook computed at
$Z/Z_\odot\approx0.231$ (the pre-restart clump matrices and the e10
check) was therefore run at the wrong reference temperature;
see Section~\ref{sec:t12-log10-bug} for the fix and its measured
effect.

\subsection{Persistence}

A tagged particle stays flagged (and re-floored, per
Eq.~\ref{eq:floor}, every step) until
\texttt{end\_time = tag\_time\allowbreak{} + HII\_region\_rebuild\_time},
at which
point it is unflagged and becomes eligible to compete for budget again
on the star's next rebuild pass. Each star's rebuild is scheduled
independently off its own age and its own last-rebuild time
(\code{feedback\_common.c}); there is no cross-star coordination or
deliberate synchronous/asynchronous toggle like the one
\citet{Smith2021} implemented (they found no difference between the
two). For the co-located, simultaneously-born stars in the validation
example this produces synchronous rebuilds in practice, simply because
the stars are identical, not because the code enforces it.

There is no carryover of unspent budget or accumulated ``credit''
between rebuild cycles -- each rebuild resets to the star's full,
instantaneous $\dot N_{\rm ion}$. This is an equilibrium-rate model,
not a time-integrated one.

\textbf{Interior-recombination-accounting gap: assessed, not an issue
(2026-07-16).} Once a particle is ionized it never reappears as a
search candidate (it fails the ``cold'' filter in
\code{feedback\_part\_can\_be\_ionized}) until its tag expires, so its
ongoing maintenance recombination cost is never re-charged against a
star's budget at any rebuild after the one that first tagged it. This
matches FIRE's own implementation (not a deviation to fix), and is
bounded by design here: the rebuild cadence is set to track the
recombination/cooling timescale, so interior gas has no time to
meaningfully recombine away between rebuilds -- no example or test run
has shown this gap actually matters in practice. Possible refinement,
not confirmed necessary: add a small margin to \code{end\_time} beyond
the nominal \code{tag\_time + HII\_region\_rebuild\_time}, to cover
cases where SWIFT's own timestepping forces a rebuild slightly earlier
than the nominal cadence -- untested, low priority given the above.

\subsection{Adaptive rebuild cadence (opt-in 2026-07-18, removed Phase 4c
2026-07-28)}
\label{sec:adaptive-cadence}

\textbf{Removed (Phase 4c, 2026-07-28).} The count-budget redesign
(Section~\ref{sec:husmith-current} of \code{04\_validation.tex}) made
rebuild cadence a purely numerical parameter: a six-cadence sweep found
the converged answer independent of it, so this mode's correctness
motivation was gone, leaving only its per-star timestep-pinning cost.
\code{GEARFeedback:HII\_adaptive\_rebuild\_cadence} and
\code{HII\_rebuild\_safety\_factor} no longer exist;
\code{HII\_rebuild\_floor\_Myr} is kept (it floors the photon-budget
interval in fixed-cadence mode too). The description below is kept as
a historical record of the mechanism that was removed.

The fixed \code{HII\_region\_rebuild\_time\_Myr} above is a single,
problem-specific value with no principled way to choose it -- this
project has used both $0.01$ and $0.5\,$Myr across different examples.
\code{GEARFeedback:HII\_adaptive\_rebuild\_cadence=1} replaces it with a
physically-derived, per-star cadence,
\begin{equation}
  \Delta t_{\rm rebuild} = f_{\rm safety}\times
  \min(t_{\rm rec,\,pixel\text{-}min},\ t_{\rm cross}),
\end{equation}
where $t_{\rm rec}=1/(\alpha_B(T_i)\,n_H)$ (recombination timescale,
computed \emph{per active angular pixel} and combined via a minimum,
since a dense clump-facing pixel and a diffuse one are physically
different environments -- diluting via a star-wide mean would
reintroduce a milder version of the mass-biasing problem angular
splitting exists to solve), $t_{\rm cross}=R_{\rm hii}/c_s(T_i)$
(sound-crossing timescale, star-wide, since $h_{\rm hii}$ is not
tracked per-pixel), and $f_{\rm safety}$ (default $0.2$) is a
dimensionless safety factor, the direct analogue of a CFL number. The
density sample for $t_{\rm rec}$ is drawn only from not-yet-tagged
front candidates already gathered by the ordinary search pass (not from
gas already tagged by this star in an earlier rebuild) -- simpler, no
changes needed to \code{feedback\_part\_can\_be\_ionized}/the gather
functions, and arguably the physically right population for
recombination risk anyway (front gas is the marginal population, not
the already-robustly-ionized interior). A backstop floor,
\code{GEARFeedback:HII\_region\_rebuild\_floor\_Myr} (default
$10^{-4}\,$Myr), prevents a pathologically short computed interval;
this is a distinct quantity from \code{Stars:min\_star\_timestep\_Myr}
(Section~\ref{sec:t16-timestep-floor}), the unrelated numerical
step-size safety floor.

Implementation: \code{radiation\_get\_T\_collisional\_K()} (extracted
helper) and
\code{radiation\_compute\_and\_cache\_HII\_rebuild\_interval()} in
\code{radiation.c}, dispatched through a thin wrapper in
\code{feedback\_common.\{h,c\}} so \code{runner\_radiation\_feedback.c}
-- generic cell-traversal code -- never includes GEAR-internal
\code{radiation.h} directly. The non-adaptive code path is untouched
(only new \code{if}/\code{else if} branches added ahead of it), so
every already-validated result stays reproducible with the feature off
(the default).

\textbf{Verification.} A genuine \code{dt=0} deadlock was found and
fixed during verification: the cached rebuild deadline is only
refreshed by an actual rebuild pass, and if that gate stays false for a
reason unrelated to timing (photon budget exhausted, aged past
\code{HII\_region\_max\_age}) the deadline goes stale and the timestep
constraint returned exactly $0.0$ forever -- fixed by flooring the
fallback instead of returning a literal $0.0$ (Section~\ref{sec:t16-timestep-floor}
generalized this floor further). With that fixed: \code{StromgrenSphere}
and all six \code{StromgrenSphereClumpSmith2021} sweep configurations
completed cleanly with sane, bounded $r_{\rm HII}$; the e10 task-graph
diagnostic showed a real, quantified improvement (time to reach
$R_{\rm st}$ roughly $47\%$ faster) without fully closing that
diagnostic's own shape mismatch -- consistent with the search-radius
cap being a separate, also-contributing factor there.

\textbf{Not yet done (2026-07-20):} the user has indicated this should
eventually become the default (replacing the fixed cadence), but wants
more testing first -- it remains opt-in, off by default, until that
testing happens.
